\subsection{Sprint 1}

Na Sprint 1, apresentamos à stakeholder os mockups finalizados pelos \ac{ages} II, com um
Figma completo, incluindo definição de cores e tipografia. Durante a apresentação, a
stakeholder identificou que o logotipo utilizado era a versão antiga com cores similares,
mas divergentes da identidade visual atual. Foram solicitadas novas funcionalidades dentro
do escopo da Escola de Raquetes, todas aceitas pelo time por estarem alinhadas ao que
seria desenvolvido.

Na aula seguinte, realizamos uma reunião interna para identificar pontos positivos a manter:
a comunicação do time se destacou positivamente. Os colegas com mais experiência apoiaram
os demais constantemente, e o grupo do WhatsApp funcionou como canal de troca de
conhecimento, com vídeos tutoriais e referências sobre as tecnologias utilizadas.

Definiu-se a divisão das tarefas entre squads equilibrados, considerando o perfil de
interesse (frontend ou backend) e a experiência de cada \ac{ages} I. Fui alocado no squad
denominado \textit{``Laura e seus Amigos''}, com uma equipe bem distribuída e membros
experientes.

Meu grupo ficou responsável pelas telas de \textbf{listar, criar e editar professor}.
Entrei em contato com minha colega de \ac{ages} III pedindo orientação. Ela indicou um vídeo
de Angular e propôs que fizéssemos juntos em uma chamada. No dia 25/03, realizamos a
chamada com Laura e André, onde implementamos a tela de criação de professor --- uma
experiência muito enriquecedora. Na aula seguinte, criamos a tela de edição. A listagem
foi desenvolvida pelo André, que encontrou dificuldades com os ícones por não haver uma
biblioteca funcional disponível no padrão do projeto.

Na semana seguinte, ao iniciarem os \textit{pushes} no repositório, a colega Sofia (\ac{ages} IV)
identificou que as telas não seguiam o padrão do Bootstrap 5. Ela criou uma tela modelo
e demonstrou na prática como realizar as correções. Essa refatoração causou um atraso de
aproximadamente uma semana na entrega do frontend.

Em paralelo, o backend sofreu atraso pois o banco de dados não estava pronto com uma semana
para o fim da sprint. Os colegas de \ac{ages} III realizaram uma chamada e conseguiram finalizar
o banco a tempo. Com isso, o backend foi desenvolvido no final de semana seguinte, concluindo
todos os CRUDs necessários para a interação entre frontend e backend. Contudo, a integração
entre as camadas não foi concluída até o fim da sprint, ficando como débito técnico para a
próxima, em razão do atraso na entrega do banco de dados.